%%% Основні відомості для заповнення%%%

\newcommand{\theWordDone}{Виконав}               % Виконав vs Виконала

\newcommand{\theStudent}{студент}                % студент vs студентка
\newcommand{\theSYear}{4}                        % Курс: 4
\newcommand{\theGroup}{ІО-21}                    % Група
\newcommand{\theStampCode}{ІАЛЦ.467200}          % Код у штампі: код КПІ та код пристрою 



\newcommand{\thesisAuthor}{Безщасний Роман Русланович}             % Автор дипломного проекту
\newcommand{\thesisAuthorTask}{Безщасному Роману Руслановичу}       % Кому видане завдання для дипломного проектування
\newcommand{\thesisAuthorInitials}{Безщасний Р.Р.}              % Прізвище та ініціали автора дипломного проекту

\newcommand{\taskDate}{10~травня 2026 року}      % Дата видачі завдання
\newcommand{\completeDate}{  ~червня \the\year\ року}  % Дата здачі закінченого проекту
\newcommand{\approvalDate}{  ~червня 2026 року}  % Дата затвердення наказом по університету
\newcommand{\approvalNum}{}   % Номер наказу затвердення теми

\newcommand{\thisisHD}           % Завідувач кафедри:  ім’я, прізвище
{Артем ВОЛОКИТА}
\newcommand{\thisisHDInitials}   % Завідувач кафедри:  прізвище, ініціали
{Волокита~А.~М.}

\newcommand{\thesisTitle}        % Тема дипломного проєкту
{Система просторового аналізу гри в волейбол з використанням машинного навчання}
%{Інтерфейс користувача для класифікації рукописних цифр}

\newcommand{\thesisOsvitnyaPrograma}     % Освітня програма навчання
{Комп'ютерні системи та мережі}

\newcommand{\thesisSpecialtyNumber}    % Номер спеціальності
{123}

\newcommand{\thesisSpecialtyTitle}     % Назва спеціальності
{Комп’ютерна інженерія}

\newcommand{\thesisDegree}             % Освiтньо-квалiфiкацiйний рiвень
{Бакалавр}


\newcommand{\thesisOrganization}       % Назва навчального закладу
{Національний технічний університет України}
\newcommand{\thesisOOrganization} 
{«Київський політехнічний інститут імені Ігоря Сікорського»}

\newcommand{\thesisInstitute}           % Факультет
{Факультет інформатики та обчислювальної техніки}

\newcommand{\thesisDepartment}          % Кафедра
{Кафедра обчислювальної техніки}


\newcommand{\supervisorFio}            % Керівник проекту: прізвище, ім’я, по батькові
{Морозов-Леонов Олександр Святославович}
\newcommand{\supervisorInitials}       % Керівник проекту:  прізвище, ініціали
{Морозов-Леонов}
\newcommand{\supervisorRegalia}        % Керівник проекту: посада, науковий ступінь, вчене звання
{аспірант}
%TODO: перевірити всі імена
\newcommand{\visorFio}{}                 % Рецензент проекту: прізвище, ім’я, по батькові
%{Писаренко Андрій Володимирович}
\newcommand{\visorRegalia}{}             % Рецензент проекту: посада, науковий ступінь, вчене звання
%{кандидат наук, доцент}

\newcommand{\consultantFio}{}           % Консультант проекту: прізвище, ім’я, по батькові
%{Сімоненко Валерій Павлович}
\newcommand{\consultantRegalia}{}      % Консультант проекту: посада, науковий ступінь, вчене звання
%{доктор технічних наук, професор}
\newcommand{\consultantInitials}{}     % Консультант проекту: прізвище та ініціали
%{Сімоненко~В.~П.}
\newcommand{\consultantPart}{}      % Консультант проекту: назва розділу
%{Нормоконтроль}
                     

%-------- Вихiднi данi до проєкту -------------
\newcommand{\dataInput}
{
технічна документація та статистичні дані.
}

%-------- Змiст розрахунково-пояснювальної записки -------------
\newcommand{\explanatoryNoteContents}
{

\underline{Розділ 1. Аналіз теорії та обґрунтування теми}

\underline{Розділ 2. Математичне та алгоритмічне забезпечення аналітичної системи}

\underline{Розділ 3. Програмна реалізація системи}

\underline{Розділ 4. Аналіз результатів та експериментальні дослідження}
}

%-------- Перелік графічного матеріалу -------------
\newcommand{\graphicMaterial}
{
принципова, структурна та  функціональна схеми системи.
}


%-------- Консультанти розділів проєкту -------------
\newcommand{\consultantsTab}
{
\begin{table}[H]
\begin{center}
\begin{tabular}{ |>{\centering}m{0.5cm}|m{5cm}|>{\centering}m{5.5cm}|>{\centering}m{2cm}|>{\centering}m{2cm}|}
\hline
\rowcolor{lightgray!20}
  &  &  & \multicolumn{2}{c |}{Підпис, дата}  \tabularnewline \cline{4-5}
  
\rowcolor{lightgray!20} 
\multirow{-1}{*}{№} & \multicolumn{1}{c|}{\multirow{-1}{*}{\raggedright Розділ}} &\multicolumn{1}{c|}{\multirow{-1}{*}{\raggedright Консультант}}  & завдання видав & завдання прийняв \tabularnewline
 
\rowcolor{lightgray!20}  
&  &   \multicolumn{1}{c|}{\multirow{-2}{*}{\raggedright\scriptsize(прізвище та ініціали)}}  &  &    \tabularnewline
\hline
   
\rownumber & \consultantPart  & \consultantInitials &  &  \TBstrut \tabularnewline  
\hline

%\rownumber & Назва розділу консультацій   & \consultantInitials &  &  \TBstrut \tabularnewline  % Поле для іншого консультанта
%\hline

\end{tabular}
\end{center}
\end{table}
}


%-------- Календарний план ------------- TODO: Потрібно заповнити календарний план(тільки останні 2 пункти)
\newcommand{\calendarPlan}
{
\begin{table}[H]
\setcounter{magicrownumbers}{0}
\begin{center}
\begin{tabular}{|>{\centering}m{0.5cm}|>{\raggedright}m{8cm}|>{\centering}m{5cm}|>{\centering}m{2cm}|}
\hline
\rowcolor{lightgray!20} № & Назва етапу виконання дипломного проекту & Термін виконання & Примітка  \tabularnewline
\hline

\rownumber & Затвердження теми проєкту  & 06.04.2026--07.06.2026 &  \TBstrut \tabularnewline 
\hline

\rownumber & Аналіз предметної області  & 01.03.2026--30.04.2026 &  \TBstrut \tabularnewline
\hline

\rownumber & Розробка архітектури та загальної структури системи  & 10.04.2026--10.05.2026 &  \TBstrut \tabularnewline
\hline

\rownumber & Програмна реалізація системи   & 20.04.2026 - 25.05.2026 &  \TBstrut \tabularnewline
\hline

\rownumber & Оформлення пояснювальної записки  & 18.05.2026-30.05.2026 &  \TBstrut \tabularnewline
\hline

\rownumber & Передзахист  &  &  \TBstrut \tabularnewline 
\hline

\rownumber & Захист  &  &  \TBstrut \tabularnewline 
\hline

%\rownumber &   &  &  \TBstrut \tabularnewline  % поле для нового етапу
%\hline
\end{tabular}
\end{center}
\end{table}
}